\chapter{Information-Theoretic Perspective}
\label{chap:information-theory}

\section{Emergent Gravity from Information Theory}
\label{sec:emergent-gravity-info}

Link to Verlinde's entropic gravity programme and related work on emergent gravity from an information-theoretic perspective. Connection to the entropic-gravity signatures discussed in \hyperref[sec:entropic-gravity]{Section}.

\section{Local Planck Length as Minimal Information Unit}
\label{sec:planck-info-unit}

Local Planck length $\ell_P(p)$ as the minimal quantisable size of an information unit at a point $p \in M$. Relation to the substrate parameter: $\ell_P \sim \hbar_q$.

\section{Time as Information Transfer}
\label{sec:time-info}

Time as a measure of information transfer and change. Time slows at heat death ($q \to 1$) because information transfer rate decreases. Time runs fast at the Big Bang ($q \to \infty$, equivalently $\hbar_q \to 1$) because information transfer is high; torsion of $E$ provides the driving mechanism.

\section{Big Bang as Initial Information State}
\label{sec:big-bang-info}

At $t = 0$: no information is known about the state of the QOU; matter is concentrated with high torsion pressure. At $t > 0$: the bundle $E$ acts on fermions, information begins to propagate, and the brane expands.

\section{Planck Units and Computational Mechanics}
\label{sec:planck-computation}

Develop the analogue of the Planck length for time and clarify its role in computational mechanics.
